\scanpage{77}
% archive sheet number: 78
/ op. hermitiens positifs, enfin satisfait à la formule (12).

\margin{Mettre en tête 2 et corollaire}

\uline{Proposition 3} Soient $E$ et $F$ deux espaces de Hilbert, soient $u_i$ ($i=1,\ldots n$) des applications de Fredholm de $E$ dans $F$\unsure{;|,} alors $u_1\wedge\ldots\wedge u_n$ est une application de Fredholm de $\overset{n}{\Lambda}^{(2)}E$ dans $\overset{n}{\Lambda}^{(2)}F$, et

\[
\lVert u_1\wedge\ldots\wedge u_n\rVert_1\leq\frac1{n!}\lVert u_1\rVert_1\ldots\lVert u_n\rVert_1\tag*{(13)}
\]

En effet, il suffit de \deleted{\unsure{$\leq$}} vérifier \deleted{que si} si \deleted{$u_1=a_1\otimes$} $u_i=\bar b_i\otimes a_i$, \deleted{\unsure{les } \illegible{1}} avec $\lVert\bar b_i\rVert\leq1$, $\lVert a_i\rVert\leq1$, alors $\lVert u_1\wedge\ldots\wedgeu_n\rVert_1\leq1/n!$, \deleted{\unsure{car}} puisqu'alors il existe une application \unsure{de Banach|de forme} \unsure{linéaire|n-linéaire} \deleted{\unsure{continue}} \inserted{de norme $\leq1/n!$} de $(\bar E\widehat{\otimes}E)^n$ dans \deleted{\unsure{dans}} $\bar F\widehat{\otimes}F$ (où $F=\overset{n}{\Lambda}^{(2)}E$) qui coïncide avec \inserted{$(u_1,\ldots u_n)\longrightarrow$} $u_1\wedge\ldots\wedge u_n$ pour des $u_i$ décomposables, donc partout puisque les deux applications sont continues \deleted{\unsure{par}} en \deleted{\unsure{tant}} qu'applications dans $L(F)$.

Or pour $u_i=\bar b_i\otimes a_i$, on a \deleted{\unsure{$\overset{n}{\Lambda}$}}

\[
u_1\wedge\ldots\wedge u_n=\frac1{n!}(\bar b_1\wedge\ldots\wedge\bar b_n)\otimes(a_1\wedge\ldots\wedge a_n),
\]

\inserted{(voir Chap I, N\textsuperscript{o} 1, formule (8))} d'où

\[
\lVert u_1\wedge\ldots\wedge u_n\rVert_1\leq\frac1{n!}\lVert\bar b_1^{\unsure{i|c|'|{}}}\wedge\ldots\wedge\bar b_n\rVert\,\lVert a_1\wedge\ldots\wedge a_n\rVert\leq\frac1{n!}\prod\lVert a_i\rVert\,\lVert\bar b_i\rVert
\]

(formule (\unsure{5,|5}),), d'où (13).

\uline{Corollaire 1} Soit $E$ un espace de Hilbert, et \deleted{\unsure{$u$}}\inserted{$u,u_i$} $\in L^1(E)$ \deleted{\unsure{de Fredholm}} $(i=1,\ldots n)$ \deleted{des} opérateurs de Fredholm dans $E$. Alors $\overset{n}{\Lambda}u$ est un opérateur de Fredholm dans $\overset{n}{\Lambda}^{(2)}E$, et on a

\[
\deleted{\unsure{\text{alors}}}\ \deleted{\unsure{\alpha_n(u)}}\left\{\begin{aligned}
\alpha_n(u_1,\ldots,u_n)&=\operatorname{Tr}(u_1\wedge\ldots\wedge u_n)\\
\alpha_n(u)&=\operatorname{Tr}(\overset{n}{\Lambda}u)
\end{aligned}\right.\quad\text{en particulier}\tag*{(14)}
\]

(où $\alpha_n(u)$ est défini dans Chap 1, \S~2, prop 1) En effet, (14) est vraie \deleted{pour \unsure{tout}} quand on
